% Chapter 5 -- Discussion and Conclusion (\input by main.tex): discussion + limitations
% + RQ answers + contributions + future work (conclusion half merged in below,
% formerly sec_9_conclusion.tex). Synthesises the detection and completion results into
% a single statement of what the thesis supports, then states every audit-identified
% limitation BY NAME so the committee cannot surface one the thesis has not already named.
%
% Honesty-phrasing rules applied (personal_agenda P1):
%   - claims matched to evidence tier (measured / inferred / speculative) explicitly
%   - "mean IoU" always qualified as "of matched detections"
%   - recall denominator restated (>=10 surviving points, Sec 4.1) wherever recall is cited
%   - completion claims scoped per the T9c outcome (PARTIALLY HOLDS) and to static cars
%   - negative results (#40/#45 etc.) presented as findings, not apologised for; #37 as a
%     metric-design lesson
%   - every number traces to a finding / artifact (cited Finding~\#NN)
%   - the six mandatory named limitations are all present: detection selection-exposure
%     (now bounded by LOSO cross-validation, Sec 4.1: precision-only, ~2 pts), statics-only
%     accuracy, compact-band d2, empty long band, 23% length-estimate fallback (#41),
%     SemanticKITTI label-quality assumption -- plus the clustering-objectness ceiling
%     (recall analysis, Ch 4) and the offline-runtime scope
%
\chapter{Conclusion and Future Work}
\label{ch:conclusion}

The preceding chapters reported results; this one interprets them and closes the
thesis. It first draws the conclusions, answering the two research questions and
restating the contributions (Section~\ref{sec:conclusion-rqs}); then gathers what the
evidence supports and how far each claim reaches, together with every limitation the
work carries, into a single section (Section~\ref{sec:disc-limitations}) --- the
validity of the evaluation, the scope of the completion claims, the ground-truth
assumptions, and the architectural and operational limits --- and finally turns the
most consequential boundaries into future work (Section~\ref{sec:conclusion-future}).
Every weakness named here is measured elsewhere in the document and cross-referenced to
the finding that quantifies it: a limitation the thesis names, quantifies, and bounds is
defended; one a committee surfaces first is not.

\section{Conclusion}
\label{sec:conclusion-rqs}

Chapter~\ref{ch:introduction} posed two research questions, both about measuring completion
on real LiDAR; each is answered in full in Section~\ref{sec:donor-validation} and restated
here in the form asked, with its reach discussed in Section~\ref{sec:disc-what-holds}.

\paragraph{RQ1 --- how to evaluate completion on real LiDAR.}
Reference-based Chamfer distance against accumulated pseudo-ground-truth is invalid, because
the accumulated reference is occluded on the same side as the input and so rewards
under-completion; the thesis replaces it with the donor-frame occluded-side coverage metric,
which scores recovery of surface seen only in other frames of the same static car and passes
a four-item validation gate.

\paragraph{RQ2 --- whether completion recovers unseen surface, and whether the benefit holds out.}
Yes, within a named scope. On sequence~08 the donor metric shows PCN recovers genuinely
occluded surface (coverage at \SI{0.1}{\metre} \num{0.000}~$\rightarrow$~\num{0.304} against
the mirror baseline \num{0.043}) and improves the amodal box
(\num{0.725}~$\rightarrow$~\num{0.771}); under a pre-registered test on held-out sequence~00
the benefit \textsc{partially holds} --- both metrics replicate, but the long band was
untestable and the compact guard fails. Moving cars are a plausibility check only.

\paragraph{Diagnostic findings.}
Two findings emerged along the way (Section~\ref{sec:recall-analysis},
Chapter~\ref{ch:evaluation}): recall is clustering-bound, not classifier-bound --- HDBSCAN
splits \SIrange{31}{37}{\percent} of cars and a chain of repair strategies fails to
generalise --- and synthetic pretraining is redundant once real clusters are plentiful, the
sim-to-real gap being total and symmetric.

\paragraph{Contributions.}
The thesis makes a single scientific contribution, methodological rather than architectural:
\textbf{a validated metric for completion quality on real LiDAR}, through which learned
completion is shown to recover genuinely occluded surface and improve amodal boxes on static
cars, partially replicated on a held-out sequence. Two diagnostic findings sit alongside it
--- the clustering-bound recall account and the redundancy of synthetic pretraining --- with
engineering contributions (the modular pipeline and its runtime characterisation, the amodal
box builder, the reproducibility infrastructure) and documented negative-result lines.

\section{Limitations and Discussion}
\label{sec:disc-limitations}

\subsection{Established Results and Their Reach}
\label{sec:disc-what-holds}

Two threads carry the strongest evidence, both belonging to the metric contribution. The
first is methodological: the donor-frame occluded-side coverage metric
(Section~\ref{sec:donor-validation}). We are not aware of a prior metric that scores a real
LiDAR car's completion against donor frames of the same static instance --- the reference-free
measures reviewed in Section~\ref{sec:bg-completion} all score against the observed input, a
synthetic prior, or the model's own outputs, and the closest prior construction (Ren et
al.~\cite{ren2022selfsupervised}) scores the accumulated cloud as a whole, not the
never-observed surface. It also carries a measured negative result about metrics: the obvious
alternative --- Chamfer distance against accumulated pseudo-ground-truth --- is invalid,
because accumulated LiDAR is itself one-sided and rewards under-completion (the raw partial
scored the lowest Chamfer distance on every real example), as do the field's standard
input-referenced UCD/UHD (Section~\ref{sec:reffree-compare}). The second thread is empirical
and rests on that metric: completion adds genuinely unseen surface (median occluded-side
coverage within \SI{0.1}{\metre} of \num{0.304} on sequence~08, against \num{0.043} for the
symmetry mirror and \num{0.000} for the raw partial) and improves the amodal box (median BEV
IoU \num{0.725}~$\rightarrow$~\num{0.771}, gains concentrated where the sensor saw least);
both survive a pre-registered held-out replication with named caveats (\textsc{partially
holds}; Section~\ref{sec:disc-eval-validity}).

The detection pipeline is a strong engineering result rather than a claim of
state-of-the-art accuracy: the supervised networks positioned in
Section~\ref{sec:bg-positioning} reach higher accuracy at the dense annotation cost this
pipeline avoids. It reaches precision \num{0.905}, recall \num{0.730}, $F_1$~\num{0.808} on
the full sequence~08 (mean IoU of matched detections \num{0.912}; recall denominator: car
instances with at least \num{10} points surviving preprocessing,
Section~\ref{sec:eval-protocol}), with an interpretable, near-annotation-free clustering
front end on modest hardware. Its recall shortfall is root-caused, not mysterious --- HDBSCAN
splitting large, close cars into fragments (Section~\ref{sec:recall-analysis}) --- and that
diagnosis, with the finding that synthetic pretraining is redundant at scale, are the two
secondary results the investigation produced.

These reaches vary. The detection operating point is the broadest, cross-validated across
all eleven labelled sequences (Section~\ref{sec:eval-loso}). Two completion results transfer
by argument rather than by sample --- the invalidity of accumulated-Chamfer follows from the
one-sidedness of accumulated LiDAR, and the donor metric is dataset-agnostic wherever static
cars are seen from multiple frames --- whereas the recall diagnosis is an inductive result
from two sequences, mechanistic but not a proven invariant. The tuned constants (the length
prior, the estimate quantile, the geometric thresholds) are sequence-derived and should be
re-checked on a new sensor, and the completion accuracy claims are scoped to static,
gate-passed cars in the compact and normal bands. For practice, three readings follow: the
interpretable, near-annotation-free detector is a viable choice where labelling budget or
inspectability outweighs peak recall, provided the domain is dense-scene; the donor metric is
a reusable instrument for any real-LiDAR completion where static objects recur across frames;
and because input quality, not model choice, governs whether completion is usable, effort is
better spent on segmentation than on swapping architectures. These reaches are bounded by the
limitations set out next.

\subsection{Validity of the Evaluation}
\label{sec:disc-eval-validity}

Detection is cross-validated, and the selection exposure is precision-only. SemanticKITTI
labels only sequences \numrange{00}{10}, so no labelled sequence can be reserved as a
conventional held-out detection test set, and the final classifier was historically selected
on sequence~08; leave-one-sequence-out cross-validation (Section~\ref{sec:eval-loso}) bounds
the exposure, the fold-08 classifier matching its recall (\num{0.737} versus \num{0.730}) and
losing only about two precision points (\num{0.881} versus \num{0.905}), so the selection
touched precision, not the clustering-imposed recall ceiling (\num{0.775}) that accounts for
most missed cars. Completion is held out more strongly still --- its claims are re-tested on
sequence~00 under a pre-registration and the model never saw a real sequence --- and that
replication returned \textsc{partially holds}: both primary metrics were decisive wins (BEV
IoU \num{0.739}~$\rightarrow$~\num{0.766}, $p=\num{1.6e-3}$; donor coverage
\num{0.000}~$\rightarrow$~\num{0.413}), but the long band was empty (no car
$\geq$~\SI{4.6}{\metre} survived the static-car construction), leaving long-car behaviour
untestable off the tuning sequence, and the compact-band guard fails there as on sequence~08.
The completion value is thus replicated once, not across a range of sequences.

\subsection{Scope of the Completion Claims}
\label{sec:disc-completion-scope}

The quantitative completion claims are bounded in what they can score. Both metrics require a
static car by construction, so moving cars cannot be scored by either and are reported only as
a plausibility check, where they reach plausible boxes at essentially the static rate
(\SI{57.9}{\percent} versus \SI{53.7}{\percent}) --- a parity result, not an accuracy claim,
and a conservative one, since the axis-aligned flag penalises the non-axis-aligned cars movers
over-represent. The metric also scores only clean, gate-passed clusters, and whether a car
reaches completion as one is decided upstream by the clustering that caps recall: the
ablations (Section~\ref{sec:abl-completion}) show the L-shape input gate sharply raises
plausible-completion rates on every sequence whereas swapping PCN for PoinTr changes them not
at all, so widening completion's reach is a segmentation problem, not a
completion-architecture one.

Within that scope, three residuals bound accuracy. The per-car length estimate falls back to
the \SI{4.14}{\metre} population prior on tracks with fewer than five gate-passed frames ---
\SI{23.0}{\percent} of sequence-08 tracks --- reintroducing compact overshoot there; the cost
sits inside the headline gain (median BEV IoU \num{0.725}~$\rightarrow$~\num{0.771}, measured
on the same tracks) and falls on the sparsest, most distant cars where the box matters least.
Even with an unbiased estimate the completed cloud stays slightly short at the far end by an
amount that scales with car length, and a pre-registered fix removed this on normal and long
cars but over-extended compacts, so the compensation is irreducibly length-dependent and the
far end remains the weakest-covered region. Both residuals compound the model's synthetic-only
training --- the source of the evaluation's independence but also the synthetic-to-real gap
(Section~\ref{sec:bg-completion}) behind the residual flatness of real completions --- which
shows most sharply on compact cars, where completions spray points outside the box at roughly
\num{100} times the mirrored baseline under the fixed prior and about \num{16} times under the
per-car estimate; no configuration clears the compact bar, so we report the ratios rather than
tune the guard, a band-localised failure the original pooled-median guard was blind to.

\subsection{Ground-Truth and Dataset Assumptions}
\label{sec:disc-gt}

Three assumptions underlie the ground truth. The amodal box is constructed by accumulating a
static car's frames, since no external amodal reference exists (KITTI raw tracklets for the
sequence-08 drive, \texttt{2011\_09\_30\_drive\_0028}, do not exist on the official archive);
it rests on the paired design comparing raw against completed on the \emph{same} box,
viewpoint- and motion-coverage guards, and a sane dimension distribution, so the claim is
robustness to plausible reference error, not an error-free reference. Recall is reported
against eligible cars --- those with at least \num{10} points surviving preprocessing ---
which excludes \SI{25.5}{\percent} of annotated cars too sparse or distant to reach the front
end, so recall against all annotated cars is approximately \num{0.54} against the \num{0.730}
eligible figure (both reported; the two ``$\geq$10-point'' rules are distinct). And all
counts, matches, and the amodal construction assume SemanticKITTI's labels are correct; we did
not audit them, and mislabelled cars would propagate into both the detection metrics and the
ground truth.

\subsection{Architectural and Operational Limits}
\label{sec:disc-architectural}

The remaining limits are architectural or operational. The recall ceiling is architectural,
not a tuning failure (Section~\ref{sec:recall-analysis}): coarsening the cluster voxel to
\SI{0.10}{\metre} recovered part of the loss (\num{0.699}~$\rightarrow$~\num{0.730}) but
coarsening further fuses adjacent cars, and no fixed-radius density clusterer reaches HDBSCAN's
recall at any tested radius, so the ceiling (\num{0.775} before any learned stage, splitting
\SIrange{31}{37}{\percent} of cars) is the measured price of the interpretable,
near-annotation-free design the thesis exists to evaluate, and lifting it means a learned
detector --- a different system. The operational domain is dense scenes: recall falls to
\num{0.336} on the sparse highway sequence~01 against \num{0.761} on dense-urban sequence~00,
mostly because highway cars are distant and recall collapses with range (\num{0.196} beyond
\SI{40}{\metre}) and partly because the front end was tuned on dense data, though the failure
is a conservative miss, not a false alarm (sequence~01 precision \num{0.951}). The system is
offline by design, running at about \SI{623.5}{\milli\second} per frame
(Section~\ref{sec:results-runtime}) with the $\leq\SI{400}{\milli\second}$ target unreachable
through the authorised tiers, and real-time operation is out of scope. Finally, the centroid
tracker feeds the length estimate and drives the \SI{23}{\percent} fallback but is never
scored as a tracker (no MOTA or IDF1), so its contribution to that fallback is unmeasured.

% Conclusion moved to the top of this chapter (Section~\ref{sec:conclusion-rqs}) to
% mirror the gospel's Conclusion-first order; the Established Results discussion is
% folded into Limitations and Discussion above. Honesty-phrasing rules (personal_agenda
% P1) still apply to the moved Conclusion: held-out outcome PARTIALLY HOLDS (#42),
% moving cars = plausibility only (#44), negatives as findings, claims within C1-C11.

\section{Future Work}
\label{sec:conclusion-future}

The limitations above named the boundaries of the present results. Three of them are the natural next
steps, each already scoped by a finding rather than left open.

\paragraph{A real-time variant.} The pipeline runs offline at roughly
\SI{623.5}{\milli\second} per scan (Section~\ref{sec:results-runtime}), dominated by the
classical CPU stages. The direct speed path is an implementation one: the residual
$\sim$\SI{185}{\milli\second} attributed to clustering is the cost of the \emph{Python}
HDBSCAN, not the HDBSCAN \emph{algorithm}, so a native or GPU implementation would attack
that floor without changing what is computed. Swapping in a faster fixed-radius clusterer is
not a free substitution --- the benchmark shows none reaches HDBSCAN's recall at any tested
radius --- so a real-time variant must preserve the HDBSCAN algorithm via a faster
implementation, or replace the front end with a learned detector, rather than trade measured
recall for speed.

\paragraph{Real-data fine-tuning for completion.} The completion model was trained on
synthetic KITTI-like partials only, which is at once the source of the donor metric's
validity and a ceiling on real-data crispness: on the sparsest, most one-sided inputs the
completions stay flat, which traces to the remaining synthetic-to-real
domain gap. Closing it would call for real-data fine-tuning, which is itself a documented
negative-result line in this project; it is a contingency
direction, to be attempted only with the leakage controls that make real-data supervision
safe.

\paragraph{Long-band amodal ground truth on a second held-out sequence.} The held-out
replication (Section~\ref{sec:results-heldout}) partially held in part because the long length band was empty on
sequence~00 --- no car of at least \SI{4.6}{\metre} survived the well-observed-static-car
construction --- so the per-car length estimate's behaviour on long cars could not be tested
off the tuning sequence. Building amodal ground truth for a third sequence with long cars
present would close that single coverage gap and upgrade the long-band claim from a
sequence-08 result to a replicated one. It is deferred here as a bounded data-construction
task, not a change to the method.

\bigskip

None of the thesis's results rests on an additional point of metric improvement: the
methodological contribution is a validated real-data completion metric, and each empirical
claim --- that learned completion adds genuinely unseen surface and improves amodal boxes on
static cars, that recall is clustering-bound, that synthetic pretraining is redundant at
scale --- is defended by explanation and, where the design allowed it, by a pre-registered
test whose outcome was reportable either way. That is the standard the work set for itself,
and the standard on which it asks to be judged.
